% Felipe Muñoz Casasola
% This is free and unencumbered software released into the public domain.

% Anyone is free to copy, modify, publish, use, compile, sell, or
% distribute this software, either in source code form or as a compiled
% binary, for any purpose, commercial or non-commercial, and by any
% means.

% In jurisdictions that recognize copyright laws, the author or authors
% of this software dedicate any and all copyright interest in the
% software to the public domain. We make this dedication for the benefit
% of the public at large and to the detriment of our heirs and
% successors. We intend this dedication to be an overt act of
% relinquishment in perpetuity of all present and future rights to this
% software under copyright law.

% THE SOFTWARE IS PROVIDED "AS IS", WITHOUT WARRANTY OF ANY KIND,
% EXPRESS OR IMPLIED, INCLUDING BUT NOT LIMITED TO THE WARRANTIES OF
% MERCHANTABILITY, FITNESS FOR A PARTICULAR PURPOSE AND NONINFRINGEMENT.
% IN NO EVENT SHALL THE AUTHORS BE LIABLE FOR ANY CLAIM, DAMAGES OR
% OTHER LIABILITY, WHETHER IN AN ACTION OF CONTRACT, TORT OR OTHERWISE,
% ARISING FROM, OUT OF OR IN CONNECTION WITH THE SOFTWARE OR THE USE OR
% OTHER DEALINGS IN THE SOFTWARE.

% For more information, please refer to <https://unlicense.org/>

\documentclass{article}

\usepackage[spanish]{babel}
\usepackage{microtype}
\usepackage[pdfusetitle]{hyperref}

\title{Archivo de servicio de uWSGI}
\author{Felipe Muñoz Casasola}
\date{27 de septiembre de 2025; revisado el 21 de febrero de 2026}

\begin{document}

\maketitle
\clearpage

\section{Archivo de servicio del sistema}
Mi instalación de uWSGI no traía un archivo de servicio, utiliza un generador de \texttt{systemd} que en futuras versiones se va a eliminar, por eso he tenido que crear manualmente este servicio. \href{https://uwsgi.readthedocs.io/en/latest/Systemd.html}{Aquí} se encuentra la página de donde saqué la información para crear el servicio. Y está archivada \href{https://web.archive.org/web/20250927100911/https://uwsgi.readthedocs.io/en/latest/Systemd.html}{aquí}. Hay una breve modificación porque el archivo \texttt{emperor.ini} no lo tengo. La versión que tengo en \texttt{/etc/systemd/system/uwsgi.service} es la siguiente:

\begin{verbatim}
[Unit]
Description=uWSGI server

[Service]
ExecStart=/bin/uwsgi --ini /etc/uwsgi/searxng.ini
RuntimeDirectory=uwsgi
Restart=always
KillSignal=SIGQUIT
Type=notify
StandardError=syslog
NotifyAccess=all

[Install]
WantedBy=multi-user.target
\end{verbatim}

\section{Archivo \texttt{searxng.ini}}
El archivo \texttt{/etc/uwsgi/searxng.ini} estaba vacío tras la instalación que hice en su día de SearxNG, que he tenido que rellenar siguiendo las instrucciones de \href{https://docs.searxng.org/admin/installation-uwsgi.html}{la guía oficial de SearxNG}, archivado \href{https://web.archive.org/web/20250927123518/https://docs.searxng.org/admin/installation-uwsgi.html}{aquí}, con el siguiente contenido:

\begin{verbatim}
# -*- mode: conf; coding: utf-8  -*-
[uwsgi]

# uWSGI core
# ----------
#
# https://uwsgi-docs.readthedocs.io/en/latest/Options.html#uwsgi-core

# Who will run the code / Hint: in emperor-tyrant mode uid & gid setting will be
# ignored [1].  Mode emperor-tyrant is the default on fedora (/etc/uwsgi.ini).
#
# [1] https://uwsgi-docs.readthedocs.io/en/latest/Emperor.html#tyrant-mode-secure-multi-user-hosting
#
uid = searxng
gid = searxng

# set (python) default encoding UTF-8
env = LANG=C.UTF-8
env = LANGUAGE=C.UTF-8
env = LC_ALL=C.UTF-8

# chdir to specified directory before apps loading
chdir = /usr/local/searxng/searxng-src/searx

# SearXNG configuration (settings.yml)
env = SEARXNG_SETTINGS_PATH=/etc/searxng/settings.yml

# disable logging for privacy
disable-logging = true

# The right granted on the created socket
chmod-socket = 666

# Plugin to use and interpreter config
single-interpreter = true

# enable master process
master = true

# load apps in each worker instead of the master
lazy-apps = true

# load uWSGI plugins
plugin = python3,http

# By default the Python plugin does not initialize the GIL.  This means your
# app-generated threads will not run.  If you need threads, remember to enable
# them with enable-threads.  Running uWSGI in multithreading mode (with the
# threads options) will automatically enable threading support. This *strange*
# default behaviour is for performance reasons.
enable-threads = true

# Number of workers (usually CPU count)
workers = %k
threads = 4

# plugin: python
# --------------
#
# https://uwsgi-docs.readthedocs.io/en/latest/Options.html#plugin-python

# load a WSGI module
module = searx.webapp

# set PYTHONHOME/virtualenv
virtualenv = /usr/local/searxng/searx-pyenv

# add directory (or glob) to pythonpath
pythonpath = /usr/local/searxng/searxng-src


# speak to upstream
# -----------------

socket = /usr/local/searxng/run/socket
buffer-size = 8192

offload-threads = %k
\end{verbatim}

\end{document}
